% Fresh strict module checks passed; canonical integration gates remain pending.
\section{Corrected interaction-label cluster bounds and operator tensors}
\label{sec:corrected_cluster_mpo}

The interaction labels are vertices of a path. The normalized tail estimate
and the exact operator representation below concern this label cutoff.
The physical-support estimate printed in~\cite{Arad2016RigorousRG} remains
false; its counterexamples are recorded in~\cite{gap:alvv16_cluster_approximation}.

\subsection{Required labels with an empty forbidden set}

For a finite set $U$ of required interaction labels, let
\begin{align}
    R_U^H(\beta)
     & =\sum_{D\subseteq U}(-1)^{|D|}
    e^{-\beta(H-H_D)},
     & H_D                            & =\sum_{i\in D}H_i.
    \label{eq:mandatory_cluster_part}
\end{align}
This is the cluster part~(\ref{eq:hastings_cluster_part}) with the forbidden
set equal to $\varnothing$. The normalization below is an auxiliary
estimate for the interaction-label expansion related to the argument
of~\cite[lines 1270--1291 and 1935--1968]{Arad2016RigorousRG}.
It is distinct from the printed spatial approximation estimate, whose
obstructions are recorded in~\cite{gap:alvv16_cluster_approximation}.

\begin{lemma}[Normalized norm of the empty-boundary cluster part]
    \label{thm:mandatory_cluster_part_norm}
    \lean{TNLean.Lattice.ClusterExpansion.norm_clusterPart_empty_boundary_le_of_lower_bound}
    \leanok

    Let $H$ and $H_i$, $i\in U$, be Hermitian complex matrices of the same
    size. Suppose $\|H_i\|\leq1$, $H\geq\lambda I$, and $\beta\geq0$.
    Then
    \begin{align}
        e^{\beta\lambda}\|R_U^H(\beta)\|
        \leq(e^\beta-1)^{|U|}.
        \label{eq:mandatory_cluster_part_norm}
    \end{align}
    The real number $\lambda$ is any supplied lower energy. The background
    $H$ need not be the sum of the selected terms.
\end{lemma}
\begin{proof}\leanok
    \uses{def:hastings_cluster_part}
    The defining finite sum~(\ref{eq:mandatory_cluster_part}) is the
    alternating sum to which the normalized mandatory-label estimate
    applies, with no forbidden labels.
\end{proof}

\begin{lemma}[The mandatory-label part as a sum over exact supports]
    \label{thm:mandatory_cluster_part_support_sum}
    \lean{TNLean.Lattice.ClusterExpansion.clusterPart_empty_boundary_eq_sum_clusterOp}
    \leanok

    Let $K$ be a finite label set, let $H_i$, $i\in K$, be complex square
    matrices, and put $H=\sum_{i\in K}H_i$. For $U\subseteq K$ and every
    real $\beta$, the operators $E_A$ of
    Definition~\ref{def:hastings_cluster_operator} satisfy
    \begin{align}
        R_U^H(\beta)=\sum_{\substack{A\subseteq K \\U\subseteq A}}E_A.
        \label{eq:mandatory_cluster_part_support_sum}
    \end{align}
\end{lemma}
\begin{proof}\leanok
    \uses{thm:hastings_cluster_part_words, def:hastings_cluster_operator}
    At each word length, partition the words containing $U$ according to
    their exact letter set $A$. The possible sets $A$ form a finite
    family. Summing their convergent word series gives
    (\ref{eq:mandatory_cluster_part_support_sum}). Word order and repeated
    occurrences of a label are retained in this partition.
\end{proof}


\subsection{A normalized tail for an interaction-label cutoff}

The cutoff in this subsection counts interaction labels. It differs from
both physical-site support and the unrestricted estimate discussed
in~\cite{gap:alvv16_cluster_approximation}. No positivity of the individual
interactions is imposed.

\begin{definition}[Intervals of interaction labels]
    \label{def:cluster_label_intervals}
    \lean{TNLean.Lattice.ClusterExpansion.pathWindows,
        TNLean.Lattice.ClusterExpansion.mem_pathWindows,
        TNLean.Lattice.ClusterExpansion.pathIntervals,
        TNLean.Lattice.ClusterExpansion.mem_pathIntervals}
    \leanok

    Give the labels $K=\{0,\ldots,N-1\}$ their path adjacency. For an integer
    $r\geq1$, let $\mathcal W_r$ be the family of intervals of exactly $r$
    labels, and let $\mathcal I_r$ be the family of nonempty intervals of
    at least $r$ labels.
\end{definition}

\begin{lemma}[Windows and connected components on a path]
    \label{thm:cluster_label_window_geometry}
    \lean{TNLean.Lattice.ClusterExpansion.ordConnected_of_pathGraph_connected,
        TNLean.Lattice.ClusterExpansion.pathGraph_connected_of_ordConnected,
        TNLean.Lattice.ClusterExpansion.eq_Icc_of_ordConnected,
        TNLean.Lattice.ClusterExpansion.exists_interval_subset_card,
        TNLean.Lattice.ClusterExpansion.noLargeCluster_pathGraph_iff,
        TNLean.Lattice.ClusterExpansion.pathWindows_exists_marker,
        TNLean.Lattice.ClusterExpansion.card_pathIntervals_le,
        TNLean.Lattice.ClusterExpansion.clusters_biUnion_pathWindows_subset}
    \leanok

    \uses{def:cluster_label_intervals, def:hastings_cluster_of_terms}
    A nonempty support is connected in the label path if and only if it is
    an interval. Every interval with at least $r$ labels contains a member
    of $\mathcal W_r$. Hence a support has no component of size at least
    $r$ if and only if it contains no member of $\mathcal W_r$.
    The left endpoints of the members of $\mathcal W_r$ are distinct and
    belong to their respective windows. Moreover, $|\mathcal I_r|\leq N^2$.
    Every component of a union of members of $\mathcal W_r$ belongs to
    $\mathcal I_r$.
\end{lemma}
\begin{proof}\leanok
    \uses{thm:hastings_clusters_meet}
    A walk cannot cross a missing label: adjacent labels differ by one.
    Conversely, the induced path on an interval is the Hasse graph of its
    inherited order and is connected. A nonempty interval is determined
    by its minimum and maximum. Its size is the difference of these
    endpoints plus one, which gives a length-$r$ subinterval and the
    injectivity of the left endpoint for fixed length. The pair of
    endpoints gives the bound $N^2$.

    If a window meets a component of a support containing it, connectedness
    and avoidance of the component boundary force the window to lie in
    that component. This proves both the cutoff characterization and
    the assertion about unions of windows.
\end{proof}

\begin{lemma}[Inclusion and exclusion over mandatory windows]
    \label{thm:cluster_window_inclusion_exclusion}
    \lean{TNLean.Lattice.ClusterExpansion.sum_clusterOp_bad_windows_eq,
        TNLean.Lattice.ClusterExpansion.exp_sub_truncatedExp_eq_sum_clusterPart_windows,
        Finset.norm_sum_powerset_biUnion_sign_le_two_pow}
    \leanok

    \uses{thm:mandatory_cluster_part_support_sum}
    Let $K$ be finite, let $H_i$, $i\in K$, be complex square matrices,
    and put $H=\sum_iH_i$. Let $\mathcal W$ be a finite family of subsets
    of $K$. For every real $\beta$,
    \begin{align}
        \sum_{\substack{A\subseteq K                   \\
                \exists I\in\mathcal W:\ I\subseteq A}}E_A
         & =\sum_{\varnothing\ne F\subseteq\mathcal W}
        (-1)^{|F|+1}R_{\bigcup F}^H(\beta).
        \label{eq:cluster_window_inclusion_exclusion}
    \end{align}
    If the retained supports of a truncated expansion are exactly those
    avoiding every window, its error $e^{-\beta H}-M$ equals
    (\ref{eq:cluster_window_inclusion_exclusion}).
    If each member $I$ of $\mathcal W$ has a distinct selected label
    belonging to $I$, then the grouped coefficient
    \begin{align}
        c_U=\sum_{\substack{\varnothing\ne F\subseteq\mathcal W \\
                \bigcup F=U}}(-1)^{|F|+1}
        ,\qquad |c_U|\leq2^{|U|}.
        \label{eq:cluster_window_coefficient}
    \end{align}
\end{lemma}
\begin{proof}\leanok
    \uses{thm:mandatory_cluster_part_support_sum}
    Apply finite inclusion and exclusion to the events $I\subseteq A$,
    then sum over $A$. For each $F$, the intersection of these events is
    $\bigcup F\subseteq A$, whose exact-support sum is
    $R_{\bigcup F}^H(\beta)$.
    A contributing family with union $U$ consists of windows contained in
    $U$. Their selected labels are distinct elements of $U$, so there
    are at most $|U|$ such windows and at most $2^{|U|}$ contributing
    families. Each sign has absolute value one.
\end{proof}

\begin{lemma}[Weighted unions of windows]
    \label{thm:cluster_label_window_weights}
    \lean{TNLean.Lattice.ClusterExpansion.sum_pow_card_unions_pathWindows_le}
    \leanok

    \uses{thm:cluster_label_window_geometry}
    For $0\leq b\leq1$, put
    $\mathcal U_r=\{\bigcup F:\varnothing\ne F\subseteq\mathcal W_r\}$.
    Then
    \begin{align}
        \sum_{U\in\mathcal U_r}b^{|U|}
        \leq\exp(N^2b^r)-1.
        \label{eq:cluster_label_window_weights}
    \end{align}
\end{lemma}
\begin{proof}\leanok
    \uses{thm:cluster_label_window_geometry,
        thm:hastings_truncation_cluster_products}
    The component family determines $U$, its members are disjoint,
    and each belongs to $\mathcal I_r$. Thus
    \begin{align}
        \sum_{U\in\mathcal U_r}b^{|U|}
         & \leq\sum_{\varnothing\ne F\subseteq\mathcal I_r}
        \prod_{C\in F}b^{|C|}
        =\prod_{C\in\mathcal I_r}(1+b^{|C|})-1
        \leq\exp\Bigl(\sum_{C\in\mathcal I_r}b^{|C|}\Bigr)-1
        \leq\exp(N^2b^r)-1.
        \label{eq:cluster_label_weight_product}
    \end{align}
    The last inequality uses $|C|\geq r$ and
    $|\mathcal I_r|\leq N^2$.
\end{proof}

\begin{theorem}[Normalized interaction-label tail on a path]
    \label{thm:normalized_cluster_label_tail}
    \lean{TNLean.Lattice.ClusterExpansion.norm_exp_sub_truncatedExp_pathGraph_le}
    \leanok

    \uses{def:hastings_truncated_expansion, thm:mandatory_cluster_part_norm,
        thm:cluster_label_window_geometry, thm:cluster_window_inclusion_exclusion,
        thm:cluster_label_window_weights}
    Let $H=\sum_{i=0}^{N-1}H_i$, where the $H_i$ are Hermitian matrices
    with $\|H_i\|\leq1$. Suppose $H\geq\lambda I$, $r\geq1$,
    $\beta\geq0$, and $b=2(e^\beta-1)\leq1$.
    Let $M_r$ retain exactly the words whose label supports have no
    path component of size at least $r$. Then
    \begin{align}
        e^{\beta\lambda}\|e^{-\beta H}-M_r\|
        \leq\exp(N^2b^r)-1.
        \label{eq:normalized_cluster_label_tail}
    \end{align}
    The number $\lambda$ is any lower energy. No commutation of the terms
    is required for this estimate.
\end{theorem}
\begin{proof}\leanok
    \uses{thm:mandatory_cluster_part_norm, thm:cluster_label_window_geometry,
        thm:cluster_window_inclusion_exclusion, thm:cluster_label_window_weights}
    The cutoff characterization and inclusion and exclusion group the
    discarded operator by required label union.
    The left endpoints provide the distinct selected labels for the
    coefficient bound. Hence
    \begin{align}
        e^{\beta\lambda}\|e^{-\beta H}-M_r\|
         & \leq\sum_{U\in\mathcal U_r}|c_U|(e^\beta-1)^{|U|}
        \leq\sum_{U\in\mathcal U_r}b^{|U|}
        \leq\exp(N^2b^r)-1.
        \label{eq:normalized_cluster_label_tail_proof}
    \end{align}
    This changes both the normalization and the cutoff of the printed
    estimate in~\cite[lines 1270--1291 and 1935--1968]{Arad2016RigorousRG}.
    It does not assert that estimate or a matrix product representation.
\end{proof}


\subsection{Physical factors at a cut of the label path}

Let $K$ be a finite interaction-label set with graph $G$, let $H_i$ be
complex square matrices, and let $E_A$ be their exact-label operators.
For a fixed criterion of large components, write
\begin{align}
    F(B)=\sum_{\substack{A\subseteq B \\
            A\text{ has no large component}}}E_A.
    \label{eq:retained_allowed_labels}
\end{align}
Every operator in (\ref{eq:retained_allowed_labels}) is a finite sum of
the existing exact-label coefficients.

\begin{lemma}[Clusters of separated supports]
    \label{thm:clusters_separated_supports}
    \lean{TNLean.Lattice.ClusterExpansion.isCluster_union_iff_of_no_adj,
        TNLean.Lattice.ClusterExpansion.noLargeCluster_union_iff_of_no_adj,
        TNLean.Lattice.ClusterExpansion.isCluster_self_of_connected,
        TNLean.Lattice.ClusterExpansion.noLargeCluster_iff_not_large_of_connected}
    \leanok

    \uses{def:hastings_cluster_of_terms}
    Suppose no vertex of $B$ is adjacent to a vertex of $D$.
    A set $C$ is a component of $B\cup D$ if and only if it is a
    component of $B$ or of $D$. Hence $B\cup D$ has no large component
    if and only if neither $B$ nor $D$ has a large component.
    A connected nonempty set $C$ is its own component and has no
    large component if and only if $C$ itself is not large.
\end{lemma}
\begin{proof}\leanok
    \uses{thm:hastings_clusters_meet}
    The two supports have no edges between them. A connected component
    meeting one side therefore stays on that side. Conversely, a
    component of one side has no adjacent vertex in the other side.
    The assertions about a connected set follow from uniqueness of
    a component meeting it.
\end{proof}

\begin{lemma}[Finite retained sums on separated label sets]
    \label{thm:retained_separated_sum_product}
    \lean{TNLean.Lattice.ClusterExpansion.sum_retained_clusterOp_union_eq_mul}
    \leanok

    \uses{thm:clusters_separated_supports, def:hastings_cluster_operator}
    Suppose nonadjacent labels correspond to commuting matrices.
    If $B,D\subseteq K$ are disjoint and no edge joins them, then
    $F(B\cup D)=F(B)F(D)$.
\end{lemma}
\begin{proof}\leanok
    \uses{thm:clusters_separated_supports,
        thm:hastings_truncation_cluster_products}
    The maps $A\mapsto(A\cap B,A\cap D)$ and
    $(A_1,A_2)\mapsto A_1\cup A_2$ are inverse on the retained
    supports. The selector separates over the two parts, and
    $E_{A_1\cup A_2}=E_{A_1}E_{A_2}$. Distributing the two finite
    sums gives the identity.
\end{proof}

\begin{lemma}[Grouping by the component containing a label]
    \label{thm:retained_crossing_component_sum}
    \lean{TNLean.Lattice.ClusterExpansion.sum_retained_clusterOp_containing_eq}
    \leanok

    \uses{thm:clusters_separated_supports,
        thm:retained_separated_sum_product}
    Suppose nonadjacent labels correspond to commuting matrices,
    and fix $p\in K$. Let $\partial C$ be the graph boundary of $C$.
    Then
    \begin{align}
        \sum_{\substack{A\subseteq K  \\A\text{ has no large component}\\p\in A}}E_A
        =\sum_{\substack{C\subseteq K \\p\in C,\ C\text{ connected}\\C\text{ not large}}}
        E_C F\bigl(K\setminus(C\cup\partial C)\bigr).
        \label{eq:retained_crossing_component_sum}
    \end{align}
\end{lemma}
\begin{proof}\leanok
    \uses{thm:hastings_clusters_meet, thm:clusters_separated_supports}
    Choose the unique component $C$ of $A$ containing $p$ and put
    $B=A\setminus C$. Then $B$ avoids both $C$ and $\partial C$.
    Conversely, any retained $B\subseteq K\setminus(C\cup\partial C)$
    gives a retained support $C\cup B$ whose component containing $p$
    is exactly $C$. These maps are inverse, and the exact-label
    operator factors as $E_{C\cup B}=E_CE_B$.
\end{proof}

\begin{lemma}[The two exterior sides of an interval]
    \label{thm:cluster_path_exterior_sides}
    \lean{TNLean.Lattice.ClusterExpansion.pathGraph_exterior_eq_sides,
        TNLean.Lattice.ClusterExpansion.pathGraph_sides_separated,
        TNLean.Lattice.ClusterExpansion.pathGraph_compl_singleton_eq_sides,
        TNLean.Lattice.ClusterExpansion.pathGraph_point_sides_separated}
    \leanok

    \uses{thm:cluster_label_window_geometry}
    Give $K=\{0,\ldots,N-1\}$ its path adjacency. For a nonempty
    connected $C$, put
    \begin{align}
        L(C) & =\{j\in K:j+1<c\text{ for every }c\in C\},
             & R(C)                                       & =\{j\in K:c+1<j\text{ for every }c\in C\}.
        \label{eq:cluster_path_exterior_sides}
    \end{align}
    The exterior $K\setminus(C\cup\partial C)$ equals $L(C)\cup R(C)$,
    and these sets are disjoint and have no edge between them.
    Similarly, deleting a single label $p$ leaves the separated sets
    $L_p=\{j:j<p\}$ and $R_p=\{j:p<j\}$.
\end{lemma}
\begin{proof}\leanok
    \uses{thm:cluster_label_window_geometry}
    The set $C$ is an interval. A vertex outside it lies entirely
    on one side; avoiding its boundary increases the separation by
    one label. The displayed inequalities also prove disjointness
    and the absence of an edge. The assertion about a single deleted
    label follows by splitting its complement according to order.
\end{proof}

\begin{lemma}[Physical nearest-neighbour crossing decomposition]
    \label{thm:physical_cluster_crossing_decomposition}
    \lean{TNLean.Lattice.ClusterExpansion.overlapGraph_nearestNeighbor_eq_pathGraph,
        TNLean.Lattice.ClusterExpansion.commute_nearestNeighbor_of_not_adj,
        TNLean.Lattice.ClusterExpansion.truncatedExp_pathGraph_eq_crossing_components,
        TNLean.Lattice.ClusterExpansion.truncatedExp_nearestNeighbor_eq_crossing_components}
    \leanok

    \uses{thm:retained_crossing_component_sum, thm:cluster_path_exterior_sides,
        thm:retained_separated_sum_product}
    Let $H_i$, $0\leq i<N$, be operators on $N+1$ sites, with
    $H_i$ supported on the designated pair $\{i,i+1\}$. Fix an integer
    $r\geq1$ and a label $p$, corresponding to the physical cut after
    site $p$. In (\ref{eq:retained_allowed_labels}), retain supports
    whose label components have fewer than $r$ labels. For every real
    $\beta$, their actual selected expansion satisfies
    \begin{align}
        M_r=F(L_p)F(R_p)+
        \sum_{\substack{C\text{ an interval of labels} \\p\in C,\ |C|<r}}
        E_C\bigl(F(L(C))F(R(C))\bigr).
        \label{eq:physical_cluster_crossing_decomposition}
    \end{align}
    All factors are the finite sums in
    (\ref{eq:retained_allowed_labels}). Neither Hermiticity nor
    positivity of the interactions is needed for this identity.
    When $r=1$ the sum over crossing components is empty.
\end{lemma}
\begin{proof}\leanok
    \uses{thm:retained_crossing_component_sum,
        thm:cluster_path_exterior_sides, thm:retained_separated_sum_product}
    The overlap graph of the designated pairs $\{i,i+1\}$ is exactly
    the path on the interaction labels. Nonadjacent terms act on
    disjoint physical sites and commute. Split the retained supports
    according to whether they contain $p$. The first family is a
    product of the two sums on $L_p$ and $R_p$. In the second family,
    use the unique component containing $p$ and split its exterior
    into $L(C)$ and $R(C)$.
\end{proof}

\begin{lemma}[Physical support of retained side sums]
    \label{thm:retained_cluster_physical_support}
    \lean{TNLean.Lattice.ClusterExpansion.isSupportedOn_sum_retained_clusterOp,
        TNLean.Lattice.ClusterExpansion.isSupportedOn_retained_nearestNeighbor_interval,
        TNLean.Lattice.ClusterExpansion.isSupportedOn_retained_nearestNeighbor_left,
        TNLean.Lattice.ClusterExpansion.isSupportedOn_retained_nearestNeighbor_right}
    \leanok

    \uses{thm:hastings_cluster_operator_support}
    If each $H_i$ acts on a designated set $X_i$, then $F(B)$ acts
    only on $\bigcup_{i\in B}X_i$. In the nearest-neighbour case,
    suppose $a\leq i$ and $i+1<b$ for every $i\in B$. Then $F(B)$
    acts only on the physical interval $[a,b)$.
    In particular, labels $i<p$ give a factor acting only on sites
    at most $p$, and labels $p<i$ give a factor acting only on sites
    strictly greater than $p$. These conclusions apply to all the
    side factors in (\ref{eq:physical_cluster_crossing_decomposition}).
\end{lemma}
\begin{proof}\leanok
    \uses{thm:hastings_cluster_operator_support}
    Each summand $E_A$, $A\subseteq B$, is supported on
    $\bigcup_{i\in A}X_i\subseteq\bigcup_{i\in B}X_i$.
    The support algebra is closed under finite sums. For a bond
    label $i$, both endpoints $i$ and $i+1$ lie in the stated physical
    interval. The two one-sided conclusions follow from the strict
    inequalities relative to $p$.
\end{proof}

These identities use the corrected interaction-label cutoff. A retained
component of $s$ labels is designated on $s+1$ physical sites; the cutoff
is not the source's selector based on the minimal physical support of a
word product. The normalization and error estimate of
Theorem~\ref{thm:normalized_cluster_label_tail} are unchanged. The
Schmidt-rank count and a matrix product operator representation are
established next.


\begin{lemma}[Two-sided rank of a crossing interval]
    \label{thm:physical_cluster_interval_rank}
    \lean{TNLean.Lattice.ClusterExpansion.card_biUnion_nearestNeighbor_le,
        TNLean.Lattice.ClusterExpansion.operatorSchmidtRank_clusterOp_nearestNeighbor_le}
    \leanok

    \uses{thm:hastings_cluster_operator_support, thm:cluster_label_window_geometry}
    Let every site have dimension $q\geq1$. If $C$ is a connected
    interval of nearest-neighbour labels, its designated physical
    support has at most $|C|+1$ sites. If $|C|<r$, then at every
    physical cut the exact-label coefficient $E_C$ satisfies
    \begin{align}
        \operatorname{OSR}(E_C)\leq q^r.\notag
    \end{align}
\end{lemma}
\begin{proof}\leanok
    \uses{thm:hastings_cluster_operator_support, thm:cluster_label_window_geometry}
    Write $C=[a,b]$. All its bonds act within the physical interval
    $[a,b+1]$, which has $|C|+1\leq r$ sites. On a designated support
    $X$, the matrix-unit bound on the two sides of a cut $A$ gives
    \begin{align}
        \operatorname{OSR}_A(T)
        \leq\min\{q^{2|X\cap A|},q^{2|X\setminus A|}\}
        \leq q^{|X|}.\notag
    \end{align}
    The two parts partition $X$, so twice the smaller cardinality
    is at most $|X|$. Apply this to the proved support of $E_C$.
\end{proof}

\begin{lemma}[Counting crossing label intervals]
    \label{thm:physical_cluster_crossing_count}
    \lean{TNLean.Lattice.ClusterExpansion.card_crossing_intervals_le}
    \leanok

    \uses{thm:cluster_label_window_geometry}
    For a fixed label $p$, there are at most $(r-1)^2$ connected
    label sets $C$ with $p\in C$ and $|C|<r$.
\end{lemma}
\begin{proof}\leanok
    \uses{thm:cluster_label_window_geometry}
    For $r\leq1$ there are none. Otherwise send $C=[a,b]$ to
    $(p-a,b-p)$. Each coordinate belongs to $\{0,\ldots,r-2\}$,
    and the two coordinates recover the endpoints and hence $C$.
\end{proof}

\begin{lemma}[All cut ranks of the label truncation]
    \label{thm:physical_cluster_all_cut_rank}
    \lean{TNLean.Lattice.ClusterExpansion.operatorSchmidtRank_truncatedExp_nearestNeighbor_le_count,
        TNLean.Lattice.ClusterExpansion.operatorSchmidtRank_truncatedExp_nearestNeighbor_le,
        TNLean.MPS.AreaLaw.operatorSchmidtRank_cluster_truncation_all_cuts_le}
    \leanok

    \uses{thm:physical_cluster_crossing_decomposition,
        thm:retained_cluster_physical_support,
        thm:physical_cluster_interval_rank, thm:physical_cluster_crossing_count}
    Let $r\geq1$, $q\geq1$, and let the interactions have the actual
    designated nearest-neighbour supports. For every real $\beta$,
    at an interior cut the selected operator satisfies
    \begin{align}
        \operatorname{OSR}(M_r)
        \leq1+(r-1)^2q^r\leq r^2q^r.\notag
    \end{align}
    The bound $r^2q^r$ holds at every chain cut, including the two
    endpoint cuts. No Hermiticity or positivity is required.
\end{lemma}
\begin{proof}\leanok
    \uses{thm:physical_cluster_crossing_decomposition,
        thm:retained_cluster_physical_support,
        thm:physical_cluster_interval_rank, thm:physical_cluster_crossing_count}
    In the exact crossing decomposition, the product with no
    crossing component has rank at most one. Each crossing term
    has rank at most $q^r$: its central factor has this rank, and
    both exterior factors act on only one side. Subadditivity and
    submultiplicativity of operator Schmidt rank give the first
    inequality. For $r\geq1$, $1+(r-1)^2\leq r^2$ and $q^r\geq1$.
    An endpoint cut has a one-dimensional configuration space on
    one side. If $q=1$ the whole operator space is one-dimensional;
    if $r=1$ the crossing family is empty. These observations also
    cover a chain with one site and no interactions.
\end{proof}

\begin{lemma}[Operator tensors from the actual cut ranks]
    \label{thm:operator_chain_from_all_cut_ranks}
    \lean{MPOTensor.rank_pairedOperator_cutMatrix_eq,
        MPOTensor.exists_traceChainOperator_of_operatorSchmidtRank_le}
    \leanok

    Let $T$ be an operator on a chain of $L\geq1$ sites of local
    dimension $q\geq0$. Let $D\geq1$ and suppose every chain cut
    has operator Schmidt rank at most $D$. Then there exist
    site-dependent matrices $A_v^{s,t}\in\mathbb C^{D\times D}$ such that
    \begin{align}
        T_{\sigma,\tau}
        =\operatorname{tr}\bigl(A_0^{\sigma_0,\tau_0}\cdots
        A_{L-1}^{\sigma_{L-1},\tau_{L-1}}\bigr).\notag
    \end{align}
\end{lemma}
\begin{proof}\leanok
    Pair the two physical indices at each site. At every cut, the
    resulting coefficient matrix is a row and column permutation
    of the realigned operator matrix, so the two ranks agree.
    The open-chain MPS construction therefore supplies actual
    rectangular tensors of internal dimension at most $D$.
    Pad these tensors to square matrices, retaining the coefficient
    because $L\geq1$. Their scalar closing factor is multiplied
    into the first tensor. This constructs the stated tensors.
\end{proof}

\begin{theorem}[Actual matrix product operator for the corrected truncation]
    \label{thm:physical_cluster_label_mpo}
    \lean{TNLean.MPS.AreaLaw.exists_mpo_cluster_truncation,
        TNLean.MPS.AreaLaw.exists_mpo_cluster_approximation}
    \leanok

    \uses{thm:physical_cluster_all_cut_rank,
        thm:operator_chain_from_all_cut_ranks, thm:normalized_cluster_label_tail}
    Let $r,q\geq1$, and let $H=\sum_{i=0}^{N-1}H_i$ on $N+1$ sites,
    with each $H_i$ supported on $\{i,i+1\}$. For every real $\beta$,
    the actual interaction-label truncation $M_r$ has a site-dependent
    matrix product operator representation of bond dimension $r^2q^r$.
    If, in addition, the terms are Hermitian, $\|H_i\|\leq1$,
    $H\geq\lambda I$, $\beta\geq0$, and $b=2(e^\beta-1)\leq1$, the
    same represented operator satisfies
    \begin{align}
        e^{\beta\lambda}\|e^{-\beta H}-M_r\|
        \leq e^{N^2b^r}-1.\notag
    \end{align}
\end{theorem}
\begin{proof}\leanok
    \uses{thm:physical_cluster_all_cut_rank,
        thm:operator_chain_from_all_cut_ranks, thm:normalized_cluster_label_tail}
    The chain has $N+1\geq1$ sites and $r^2q^r\geq1$. Apply the
    constructed all-cut rank estimate and then the preceding tensor
    construction. Substitute the resulting exact representation in
    the already proved normalized error bound.
\end{proof}

This is a corrected interaction-label theorem related to
\cite[Theorems on cluster approximation and representation, lines 1270--1303]{Arad2016RigorousRG}.
The printed physical-support approximation remains false; its
counterexamples and the distinction between label count and minimal
physical support remain in \cite{gap:alvv16_cluster_approximation}.
The result asserts the existence of the tensors of the actual finite
selected operator. It does not assert the source's computational
running-time bound.
